\section{Second derivative and convexity}

The first derivative tells us whether a function is increasing or decreasing. The second derivative tells us how the first derivative changes. This gives information about the bending of the graph.

If the first derivative is increasing, slopes are becoming larger. If the first derivative is decreasing, slopes are becoming smaller. This is the idea behind convexity and concavity.

\subsection*{The second derivative}

If \(f'\) is differentiable, then the derivative of \(f'\) is called the second derivative of \(f\). It is written
\[
f''(x).
\]
Thus
\[
f''(x)=(f'(x))'.
\]

\begin{examplebox}
Let
\[
f(x)=x^3-3x.
\]
Then
\[
f'(x)=3x^2-3.
\]
Differentiate again:
\[
f''(x)=6x.
\]
The second derivative tells us how the slope \(f'(x)\) changes.
\end{examplebox}

\subsection*{Convexity and concavity}

If \(f''(x)>0\) on an interval, the function is convex there. The graph bends upward. If \(f''(x)<0\), the function is concave there. The graph bends downward.

\begin{theorembox}
Let \(f\) be twice differentiable on an interval.
\begin{itemize}
\item If \(f''(x)>0\) on the interval, then \(f\) is convex on that interval.
\item If \(f''(x)<0\) on the interval, then \(f\) is concave on that interval.
\end{itemize}
\end{theorembox}

This should be connected to slopes. When \(f''>0\), the slopes increase. When \(f''<0\), the slopes decrease.

\begin{examplebox}
Study the convexity of
\[
f(x)=x^3-3x.
\]
We have
\[
f''(x)=6x.
\]
Thus
\[
f''(x)<0 \quad\text{for } x<0,
\]
and
\[
f''(x)>0 \quad\text{for } x>0.
\]
Therefore the function is concave on \((-\infty,0)\) and convex on \((0,\infty)\).
\end{examplebox}

\subsection*{Inflection points}

An inflection point is a point where the graph changes convexity. Usually, this occurs where the second derivative changes sign.

\begin{definitionbox}
A point \(a\) is an inflection point of \(f\) if the graph of \(f\) changes convexity at \(a\).
\end{definitionbox}

\begin{examplebox}
For
\[
f(x)=x^3-3x,
\]
we found
\[
f''(x)=6x.
\]
The second derivative changes sign at \(x=0\). Therefore the graph has an inflection point at \(x=0\). The corresponding point on the graph is
\[
(0,f(0))=(0,0).
\]
\end{examplebox}

The condition \(f''(a)=0\) by itself is not enough. The second derivative must change sign for convexity to change.

\subsection*{Second derivative test for extrema}

The second derivative can sometimes classify critical points.

\begin{theorembox}
Suppose \(f'(a)=0\).
\begin{itemize}
\item If \(f''(a)>0\), then \(f\) has a local minimum at \(a\).
\item If \(f''(a)<0\), then \(f\) has a local maximum at \(a\).
\item If \(f''(a)=0\), the test gives no conclusion.
\end{itemize}
\end{theorembox}

The test matches the geometric picture. If the graph is convex at a horizontal tangent, it bends upward and forms a local minimum. If it is concave at a horizontal tangent, it bends downward and forms a local maximum.

\begin{examplebox}
Let
\[
f(x)=x^2.
\]
Then
\[
f'(x)=2x,
\qquad
f''(x)=2.
\]
The critical point is \(x=0\). Since
\[
f''(0)=2>0,
\]
there is a local minimum at \(x=0\).
\end{examplebox}

\begin{summarybox}
When using the second derivative, ask:
\begin{itemize}
\item What is \(f''(x)\)?
\item Where is \(f''(x)>0\) or \(f''(x)<0\)?
\item Does the graph change convexity?
\item Are there inflection points?
\item Can the second derivative test classify a critical point?
\item If \(f''(a)=0\), is another test needed?
\end{itemize}
\end{summarybox}

The second derivative helps us read the bending of a graph, not only its direction of movement.